\documentclass[12pt,a4paper]{article}
\usepackage[utf8]{inputenc}
\usepackage[T1]{fontenc}
\usepackage[brazil]{babel}
\usepackage{amsmath,amssymb}
\usepackage{geometry}
\geometry{margin=2.5cm}
\title{Material de Estudo: Apresenta\c{c}\~ao de Trabalhos I --- Estrutura e Conte\'udo}
\author{Princ\'ipio de Modelagem Matem\'atica}
\date{Unidade 32}
\begin{document}
\maketitle

\section*{Objetivo}
Aprender a estruturar uma apresenta\c{c}\~ao cient\'ifica de modelagem matem\'atica com clareza e rigor.

\section*{Estrutura Cient\'ifica de uma Apresenta\c{c}\~ao}

Uma boa apresenta\c{c}\~ao segue o mesmo padr\~ao de um artigo cient\'ifico:

\begin{enumerate}
  \item \textbf{Introdu\c{c}\~ao (15\%):} Motiva\c{c}\~ao, pergunta de pesquisa, objetivos, estrutura da apresenta\c{c}\~ao.
  \item \textbf{Materiais e M\'etodos (30\%):} Descri\c{c}\~ao do modelo, hip\'oteses, par\^ametros, implementa\c{c}\~ao.
  \item \textbf{Resultados (30\%):} Gr\'aficos, tabelas, figuras com legenda. Dados sem interpreta\c{c}\~ao.
  \item \textbf{Discuss\~ao (15\%):} Interpreta\c{c}\~ao, limita\c{c}\~oes, compara\c{c}\~ao com a literatura.
  \item \textbf{Conclus\~oes (10\%):} Resposta \`a pergunta, contribui\c{c}\~oes, perspectivas.
\end{enumerate}

\section*{Boas Pr\'aticas}

\subsection*{Slides}
\begin{itemize}
  \item M\'aximo de 6--8 linhas por slide; fonte $\ge 24$ pt.
  \item Uma ideia principal por slide.
  \item Figuras com eixos rotulados, unidades e legenda.
  \item N\~ao reproduza equa\c{c}\~oes sem explic\'a-las.
\end{itemize}

\subsection*{Fala}
\begin{itemize}
  \item N\~ao leia os slides --- conte uma hist\'oria.
  \item Mantenha contato visual com a audi\^encia.
  \item Use transi\c{c}\~oes: ``agora que mostramos X, passamos para Y''.
  \item Respeite o tempo: ensaie antes.
\end{itemize}

\subsection*{Rigor Cient\'ifico}
\begin{itemize}
  \item Explicite as hip\'oteses do modelo.
  \item Mostre valida\c{c}\~ao (compara\c{c}\~ao com dados reais ou solu\c{c}\~ao anal\'itica).
  \item Discuta limitac\~oes honestamente.
  \item Cite as refer\^encias usadas.
\end{itemize}

\section*{Crit\'erios de Avalia\c{c}\~ao}

\begin{tabular}{|l|c|}
\hline
\textbf{Crit\'erio} & \textbf{Peso} \\
\hline
Clareza e organiza\c{c}\~ao & 20\% \\
Rigor cient\'ifico do modelo & 25\% \\
Qualidade dos resultados & 25\% \\
Discuss\~ao e conclus\~oes & 20\% \\
Uso do tempo e apresenta\c{c}\~ao & 10\% \\
\hline
\end{tabular}

\section*{Exerc\'icios}

\begin{enumerate}
  \item \textbf{(Estrutura)} Pegue um artigo cient\'ifico da \'area e identifique as 5 se\c{c}\~oes. Qual propor\c{c}\~ao de p\'aginas cada se\c{c}\~ao ocupa? Isso bate com a estrutura ideal acima?
  
  \item \textbf{(Slides)} Crie um esbo\c{c}o (lista de slides) para uma apresenta\c{c}\~ao de 10 min sobre o modelo SIR. Quantos slides voc\^e usaria? O que estaria em cada um?
  
  \item \textbf{(Hip\'oteses)} Para o modelo SIR cl\'assico: liste todas as hip\'oteses impl\'icitas (popula\c{c}\~ao homog\^enea, mistura perfeita, sem vitalidade\ldots). Como cada hip\'otese limita o modelo?
  
  \item \textbf{(Figuras)} Descreva como voc\^e apresentaria graficamente os resultados de uma DM de LJ: quais gr\'aficos mostraria ($g(r)$, MSD, $T(t)$, $E(t)$), em que ordem, e o que cada um mostraria.
  
  \item \textbf{(Limita\c{c}\~oes)} Quais s\~ao as 3 principais limita\c{c}\~oes de qualquer modelo visto na disciplina? Como a discuss\~ao dessas limita\c{c}\~oes fortalece (n\~ao enfraquece) a apresenta\c{c}\~ao?
  
  \item \textbf{(Pergunta de pesquisa)} Formule uma pergunta de pesquisa clara e test\'avel para um dos seguintes temas: (a) tr\'afego urbano com NaSch; (b) difus\~ao em l\'iquidos via DM; (c) espalhamento de epidemia via SIR-CA.
  
  \item \textbf{(Valida\c{c}\~ao)} Para o seu modelo, que dados reais ou resultado te\'orico voc\^e usaria para validar? O que seria um ``bom'' e um ``ruim'' resultado de valida\c{c}\~ao?
\end{enumerate}

\section*{Resumo R\'apido}
\begin{itemize}
  \item Estrutura: Intro $\to$ M\'etodos $\to$ Resultados $\to$ Discuss\~ao $\to$ Conclus\~ao.
  \item Slides claros, fala fluente, figuras rotuladas.
  \item Rigor: hip\'oteses expl\'icitas, valida\c{c}\~ao, limita\c{c}\~oes honestas.
\end{itemize}

\section*{Refer\^encias}
\begin{itemize}
  \item Mack, C.\ A.\ \textit{How to Write a Good Scientific Paper}. SPIE Press, 2018.
  \item Notas de aula da disciplina.
\end{itemize}

\end{document}
